\documentclass[preprint,12pt,authoryear]{elsarticle}

% ------------------------------------------------------------------
% Packages
% ------------------------------------------------------------------

\usepackage{amsmath}
\usepackage{amssymb}
\usepackage{amsthm}
\usepackage{mathtools}
\usepackage{booktabs}
\usepackage{graphicx}
\usepackage{microtype}
\usepackage[
    colorlinks=true,
    allcolors=blue
]{hyperref}

% ------------------------------------------------------------------
% Theorem environments
% ------------------------------------------------------------------

\newtheorem{proposition}{Proposition}

% ------------------------------------------------------------------
% Journal and bibliography
% ------------------------------------------------------------------

\journal{Economics Letters}
\biboptions{authoryear,round}

% ------------------------------------------------------------------
% Document
% ------------------------------------------------------------------

\begin{document}

\begin{frontmatter}

\title{
Concentration Is Not Scaling:
An Exact Decomposition of Aggregate Variance Elasticities
}

\author[aff1]{Matías Iglesias\corref{cor1}}
\cortext[cor1]{Corresponding author}
\ead{mniglesias@dc.uba.ar}

\address[aff1]{
Departamento de Computación, Facultad de Ciencias Exactas y Naturales,
Universidad de Buenos Aires, Buenos Aires, Argentina
}

\begin{abstract}
Granular models relate aggregate volatility to the concentration of
economic weights. This connection is exact only under restrictive
assumptions on marginal shock variances, dependence, and population
composition. Concentration is a state variable, whereas volatility
scaling is an elasticity defined along a population path. We show that
identical Herfindahl paths can generate different aggregate-variance
elasticities and derive an exact finite-interval decomposition of that
elasticity into changes in diagonal variance and signed covariance
components. The decomposition uses the logarithmic mean of total
variance and requires no logarithm of individual components. A
contribution-weighted average of component elasticities emerges only
under positivity and comparability conditions. Deterministic
constructions verify the direct and reconstructed elasticities to
machine precision. The result clarifies the boundary of
concentration-based granular reasoning and supplies an accounting
framework for micro-to-macro analyses of diversification.
\end{abstract}

\begin{keyword}
aggregate volatility
\sep concentration
\sep diversification
\sep variance decomposition
\sep covariance
\sep scaling elasticity

\medskip
\textit{JEL classification:} C02, C10, D20
\end{keyword}

\end{frontmatter}

% ------------------------------------------------------------------
% Main text
% ------------------------------------------------------------------

\section{Introduction}
\label{sec:introduction}

The granular approach restored an elementary but consequential fact to
macroeconomic aggregation: units enter an aggregate according to their
economic weights, so a small number of very large firms can prevent
idiosyncratic fluctuations from averaging out rapidly
\citep{gabaix2011}. Under mutually independent shocks with a common
variance, the variance of a weighted aggregate is proportional to the
Herfindahl index
\begin{equation}
    H(N)=\sum_{i=1}^{N}w_i(N)^2.
\end{equation}
A fat-tailed firm-size distribution can therefore generate a
concentration path that declines slowly as the population expands.

This benchmark is sometimes given a stronger interpretation: that the
concentration path identifies the population-scaling rate of aggregate
volatility. That conclusion is not general. Concentration describes the
distribution of weights at a given scale. Aggregate-volatility scaling
is an elasticity describing how the complete variance--covariance
structure changes along a specified population path. The two objects
coincide only when restrictions on marginal variances, dependence, and
composition supply the missing bridge between them.

For the weighted aggregate
\begin{equation}
    X_t(N)=\sum_{i=1}^{N}w_i(N)\varepsilon_{i,t}(N),
\end{equation}
aggregate variance is
\begin{align}
\label{eq:intro_variance}
    V(N)
    &=
    \sum_{i=1}^{N}w_i(N)^2
    \operatorname{Var}\!\left[\varepsilon_{i,t}(N)\right]
    \nonumber\\
    &\quad+
    2\sum_{i<j}w_i(N)w_j(N)
    \operatorname{Cov}\!\left[
        \varepsilon_{i,t}(N),\varepsilon_{j,t}(N)
    \right].
\end{align}
The proportionality $V(N)=\sigma^2H(N)$ follows when all shocks share a
variance $\sigma^2$ that is stable along the population path and the
off-diagonal terms vanish. These are not incidental simplifications.
Even under independence, allowing the common shock variance to vary with
population scale gives
\begin{equation}
\label{eq:benchmark_elasticity_bridge}
    -\frac{d\log V(N)}{d\log N}
    =
    -\frac{d\log H(N)}{d\log N}
    -\frac{d\log \sigma^2(N)}{d\log N}.
\end{equation}
The concentration elasticity identifies the variance elasticity only
when the second term is zero. Heterogeneous firm-level variances replace
$H(N)$ by a variance-weighted concentration object, while common shocks
and economic linkages add off-diagonal terms. Population expansion may
also alter the size, sectoral, or network composition of the aggregate.
Once these channels are present, the Herfindahl index is not a sufficient
statistic for either aggregate variance or its population elasticity.

This paper separates the granular insight from that stronger inference.
Granularity correctly establishes that aggregate variance is weighted
by firm shares and may be dominated by a small number of large units. It
does not, by itself, establish that the firm-size distribution identifies
the rate at which aggregate variance changes with population size. We
make this distinction precise in two steps.

First, we show that identical concentration paths can generate different
aggregate-volatility scaling. Equal-weight populations all satisfy
$H(N)=1/N$, but aggregate variance follows $1/N$ only under independence.
With persistent positive correlation, it instead approaches a positive
covariance floor and its scaling elasticity converges to zero. No
statistic constructed solely from the weights can distinguish the two
systems.

Second, we derive an exact finite-interval decomposition of the aggregate
variance elasticity into changes in diagonal variance and signed
covariance components. The formula does not take logarithms of individual
components and therefore remains well defined when covariance
contributions are negative. It also identifies the restrictive
positive-component case in which an aggregate elasticity may be read as
a contribution-weighted average of component elasticities.

The paper sits at the intersection of three established literatures.
Granular models study how unequal weights transmit microeconomic shocks
to aggregates; network models show that dependence and propagation can
prevent diversification; and index-decomposition and risk-allocation
methods provide exact additive accounting tools
\citep{gabaix2011,network_placeholder,lmdi_placeholder,risk_decomposition_placeholder}.
Our claim is deliberately narrower than invention of any of these
ingredients. The contribution is their combination for a distinct
micro-to-macro object: the finite-interval elasticity of aggregate
variance with respect to population scale. This formulation separates a
level statement about weighted exposure from a path statement about
diversification, retains signed covariance contributions, and yields an
identity that can be checked directly in empirical or synthetic data.

\section{Exact decomposition}
\label{sec:exact_decomposition}

Consider an aggregate composed of $Q$ contribution series,
\begin{equation}
    X_t(N)=\sum_{q=1}^{Q}c_{q,t}(N),
\end{equation}
where $N>0$ indexes population scale. Its variance is
\begin{equation}
\label{eq:variance_accounting}
    V(N)
    =
    \sum_{q=1}^{Q}\operatorname{Var}\!\left[c_{q,t}(N)\right]
    +
    2\sum_{q<r}\operatorname{Cov}\!\left[c_{q,t}(N),c_{r,t}(N)\right].
\end{equation}
Let $\mathcal{K}$ collect the diagonal terms
$C_{qq}(N)=\operatorname{Var}[c_{q,t}(N)]$ and the doubled off-diagonal
terms
$C_{qr}(N)=2\operatorname{Cov}[c_{q,t}(N),c_{r,t}(N)]$ for $q<r$. Then
\begin{equation}
\label{eq:component_sum}
    V(N)=\sum_{k\in\mathcal{K}}C_k(N).
\end{equation}
Diagonal components are non-negative, whereas covariance components may
have either sign. The component universe $\mathcal{K}$ is assumed to
retain a comparable meaning at the population levels being contrasted.

For $N_0,N_1>0$, $N_0\neq N_1$, and positive aggregate variance at both
levels, define the finite-interval variance elasticity
\begin{equation}
\label{eq:direct_alpha}
    \widehat{\alpha}(N_0,N_1)
    =
    -\frac{
        \log V(N_1)-\log V(N_0)
    }{
        \log N_1-\log N_0
    }.
\end{equation}
For $a,b>0$, let
\begin{equation}
\label{eq:log_mean}
    L(a,b)
    =
    \begin{cases}
        \displaystyle\frac{a-b}{\log a-\log b}, & a\neq b,\\[1em]
        a, & a=b,
    \end{cases}
\end{equation}
be the logarithmic mean.

\begin{proposition}[Exact finite-interval decomposition]
\label{prop:exact_decomposition}
For each $k\in\mathcal{K}$, define
\begin{equation}
\label{eq:component_alpha_contribution}
    \widehat{\alpha}^{\,C}_k(N_0,N_1)
    =
    -\frac{
        C_k(N_1)-C_k(N_0)
    }{
        L\!\left(V(N_1),V(N_0)\right)
        \left[\log N_1-\log N_0\right]
    }.
\end{equation}
Then
\begin{equation}
\label{eq:exact_alpha_decomposition}
    \widehat{\alpha}(N_0,N_1)
    =
    \sum_{k\in\mathcal{K}}
    \widehat{\alpha}^{\,C}_k(N_0,N_1).
\end{equation}
The identity remains valid when individual covariance components are
negative.
\end{proposition}

\begin{proof}
Equation~\eqref{eq:component_sum} implies
\begin{equation}
    \sum_{k\in\mathcal{K}}
    \left[C_k(N_1)-C_k(N_0)\right]
    =V(N_1)-V(N_0).
\end{equation}
Summing equation~\eqref{eq:component_alpha_contribution} over $k$ and
using
\begin{equation}
    \frac{V(N_1)-V(N_0)}{L(V(N_1),V(N_0))}
    =\log V(N_1)-\log V(N_0)
\end{equation}
produces equation~\eqref{eq:direct_alpha} exactly.
\end{proof}

The decomposition attributes the aggregate elasticity to changes in the
underlying variance and covariance components without requiring their
individual logarithms. Signed covariance terms therefore enter directly,
provided total variance remains positive. The result is exact for each
finite interval and does not require a global power law. A regression
coefficient estimated over many population levels may still provide a
useful summary, but it averages interval-specific elasticities whose
component weights can change across the fitted range. The decomposition
therefore distinguishes an accounting identity from a parametric claim:
it reconstructs what drives the observed change in aggregate variance
without assuming that one invariant exponent governs every scale. In the
differentiable limit,
\begin{equation}
\label{eq:local_signed_decomposition}
    \alpha_{\mathrm{eff}}(N)
    =-
    \frac{d\log V(N)}{d\log N}
    =-
    \sum_{k\in\mathcal{K}}
    \frac{N C_k'(N)}{V(N)}.
\end{equation}
If every component is positive, defining
$\alpha_k(N)=-d\log C_k(N)/d\log N$ gives
\begin{equation}
\label{eq:weighted_local_alpha}
    \alpha_{\mathrm{eff}}(N)
    =
    \sum_{k\in\mathcal{K}}
    \frac{C_k(N)}{V(N)}\alpha_k(N).
\end{equation}
This resembles a weighted average, but its weights generally vary with
$N$. A fixed average of local exponents arises only under additional
restrictions, such as common component elasticities or contribution
shares that remain invariant along the population path. When positive
components decay at different power-law rates, their shares change with
$N$ and the slowest-decaying component eventually dominates. With signed
covariance components the normalized contributions need not be
non-negative, so the representation is not a convex combination and the
aggregate elasticity need not lie within the range of the diagonal
elasticities. Covariance is therefore not merely a correction to the
level of aggregate variance; its own scale dependence can determine the
limiting diversification rate.

Comparability across scales is substantive. If parts appear, disappear,
or change meaning between $N_0$ and $N_1$, the measured elasticity may
combine diversification with composition, partition, or dependence
changes. The identity remains valid only after those changes are
represented through a common accounting structure.

\section{Why concentration does not identify volatility scaling}
\label{sec:concentration_scaling}

Consider the weighted aggregate
\begin{equation}
\label{eq:weighted_aggregate}
    X_t(N)
    =
    \sum_{i=1}^{N}w_i(N)\varepsilon_{i,t}(N),
    \qquad
    \sum_{i=1}^{N}w_i(N)=1.
\end{equation}
For equal weights, $w_i(N)=1/N$, the concentration path is
\begin{equation}
\label{eq:equal_weight_concentration}
    H(N)=\frac{1}{N}.
\end{equation}
Suppose first that shocks have unit variance and are mutually
independent. Aggregate variance is then
\begin{equation}
\label{eq:iid_variance}
    V_{\mathrm{iid}}(N)
    =
    \operatorname{Var}\!\left[
        \frac{1}{N}\sum_{i=1}^{N}\varepsilon_{i,t}
    \right]
    =\frac{1}{N},
\end{equation}
and its local variance elasticity is
\begin{equation}
    \alpha_{\mathrm{iid}}(N)=1.
\end{equation}

Now retain the same weights and marginal variances but impose a common
pairwise correlation $\rho\in(0,1)$. Aggregate variance becomes
\begin{align}
\label{eq:correlated_variance}
    V_{\rho}(N)
    &=
    \frac{1}{N^2}\left[N+\rho N(N-1)\right]
    \nonumber\\
    &=
    \rho+\frac{1-\rho}{N},
\end{align}
with local elasticity
\begin{equation}
\label{eq:correlated_alpha}
    \alpha_{\rho}(N)
    =
    \frac{1-\rho}{\rho N+1-\rho}.
\end{equation}
The independent aggregate continues to diversify at rate $N^{-1}$,
whereas the correlated aggregate approaches the variance floor
$\lim_{N\to\infty}V_{\rho}(N)=\rho$ and
$\lim_{N\to\infty}\alpha_{\rho}(N)=0$.

\begin{figure}[t]
    \centering
    \includegraphics[
        width=0.88\linewidth
    ]{figures/figure_1_concentration_is_not_scaling.pdf}
    \caption{
    Equal concentration does not imply equal volatility scaling.
    Both populations have concentration $H(N)=1/N$. Aggregate variance
    declines as $1/N$ under independent shocks but approaches a positive
    floor when shocks share a common correlation component.
    }
    \label{fig:concentration_scaling}
\end{figure}

The two economies have the same number of units, weights, Herfindahl
index, effective number of firms, and entire concentration path. Yet they
have different variance paths and limiting elasticities. The missing
information is not a refinement of concentration: it is the
scale-dependent variance--covariance structure of the contributions.
This distinction also rules out a common empirical shortcut. Estimating a
power law for $H(N)$ and interpreting its slope as the population exponent
of aggregate variance is valid only after the micro-variance and
covariance terms have been shown to be scale neutral. Concentration and
aggregate-volatility scaling are therefore different estimands even when
they coincide under the independent homogeneous benchmark.

This result qualifies rather than overturns the granular mechanism. The
presence of large firms may still make weighted idiosyncratic shocks
quantitatively important. What does not follow is that the population
elasticity of aggregate variance can be recovered from the firm-size
distribution alone. The Herfindahl connection is conditional on stable
homogeneous shock variances, independence, and a coherent population
path. Relaxing those conditions does not merely add noise around the
granular prediction; it changes the estimand that concentration can
identify.

\section{Deterministic computational verification}
\label{sec:computational_verification}

We verify the variance and scaling identities using deterministic
finite-sample constructions rather than Monte Carlo approximations. For
$N\in\{4,8,16,32,64\}$, we construct contribution panels with $T=256$
observations and fixed seed $20260723$. Centered random matrices are
orthogonalized by QR decomposition and scaled to have identity sample
covariance under $ddof=1$. Cholesky transformations then impose known
target covariance matrices, making measured and designed moments agree
up to floating-point precision.

The first construction reproduces the independent and correlated
systems in Figure~\ref{fig:concentration_scaling}. Two additional
constructions contain heterogeneous diagonal components, first without
and then with a signed covariance contribution. At each $N$, aggregate
variance is measured directly from the summed contribution series and
reconstructed from diagonal variances and doubled covariance terms. The
finite-interval elasticity is then computed directly from adjacent
aggregate variances and reconstructed through
Proposition~\ref{prop:exact_decomposition}.

\begin{figure}[t]
    \centering
    \begin{minipage}{0.49\linewidth}
        \centering
        \includegraphics[
            width=\linewidth
        ]{figures/figure_2_direct_vs_reconstructed_alpha.pdf}

        \textbf{(a) Direct and reconstructed elasticity}
    \end{minipage}
    \hfill
    \begin{minipage}{0.49\linewidth}
        \centering
        \includegraphics[
            width=\linewidth
        ]{figures/figure_3_alpha_component_contributions.pdf}

        \textbf{(b) Signed component contributions}
    \end{minipage}
    \caption{
    Exact reconstruction of aggregate variance scaling. Panel~(a)
    compares the directly measured finite-interval elasticity with its
    reconstruction. Panel~(b) separates diagonal and signed covariance
    contributions. Differences are at floating-point precision.
    }
    \label{fig:alpha_verification}
\end{figure}

The maximum absolute residual is $4.441\times10^{-16}$ for the variance
identity and $1.221\times10^{-15}$ for the scaling identity; the largest
discrepancy between designed and measured covariance components is
$5.551\times10^{-16}$. These calculations verify the implementation and
numerical consistency of the accounting chain. They do not constitute
empirical evidence about the scaling behavior of any particular
economy.

\section{Conclusion}
\label{sec:conclusion}

Concentration locates aggregate exposure; it does not generally identify
how aggregate variance changes as the population expands. The
Herfindahl connection requires a stable homogeneous variance structure,
restricted covariance, and a coherent population path. We provide an
exact finite-interval decomposition that attributes the aggregate
variance elasticity to changes in diagonal and signed covariance
components without logging the components themselves. The result
preserves the central granular insight while defining its boundary: the
firm-size distribution can make large firms important, but it is not a
sufficient statistic for population scaling. The decomposition supplies
a transparent accounting contract for empirical micro-to-macro work,
shifting the empirical task from reading an exponent off the size
distribution to measuring how variance and covariance contributions
change along an explicit population path.

% ------------------------------------------------------------------
% Declarations
% ------------------------------------------------------------------

\section*{Funding}

This research did not receive any specific grant from funding agencies
in the public, commercial, or not-for-profit sectors.

\section*{Declaration of competing interest}

The author declares that he has no known competing financial interests
or personal relationships that could have appeared to influence the
work reported in this paper.

\section*{Data and code availability}

No empirical data were used in this study. Source code, deterministic
computational constructions, tests, and generated verification artifacts
will be made available in a public repository.

\section*{
Declaration of generative AI and AI-assisted technologies in the
writing process
}

During the preparation of this work, the author used OpenAI ChatGPT and
Codex to assist with research organization, software implementation,
adversarial checking, and language revision. The author reviewed and
verified the mathematical derivations, computational outputs,
references, and final text, edited the resulting material as needed,
and takes full responsibility for the content of the publication.

% ------------------------------------------------------------------
% References
% ------------------------------------------------------------------

\bibliographystyle{elsarticle-harv}
\bibliography{references}

\end{document}